\documentclass[11pt]{article}
\usepackage[T1]{fontenc}
\usepackage{lmodern}
\usepackage[margin=1in]{geometry}
\usepackage{amsmath,amssymb,amsthm,booktabs,array,longtable}
\usepackage[hidelinks,breaklinks]{hyperref}
\usepackage{xurl}
\newtheorem{theorem}{Theorem}
\newtheorem{lemma}[theorem]{Lemma}
\newtheorem{corollary}[theorem]{Corollary}
\newtheorem*{consequence}{Corollary}
\newcommand{\Z}{\mathbb Z}
\newcommand{\Q}{\mathbb Q}
\newcommand{\F}{\mathbb F}
\newcommand{\coef}[2]{[x^{#1}]\,#2}
\newcommand{\ord}{\operatorname{ord}}
\newcommand{\leanroot}{https://github.com/the-omega-institute/trureturing/blob/074df09ddbbef1843ddb383577bfaa827f753e36}
\title{Hanna's OEIS A389474 FORMULA (3):\\
a modulo-nine proof and companion congruences}
\author{Haobo Ma\\\small ChronoAI Pte Ltd, Singapore
\and Wenlin Zhang\\\small National University of Singapore\\
\small ChronoAI Pte Ltd, Singapore}
\date{November 8, 2026}
\begin{document}
\maketitle
\begin{abstract}
We prove Hanna's conjecture in OEIS A389474, FORMULA (3), and the
companion conjectures in A389474, FORMULA (4),
A389475, FORMULA (4.a), (4.b), and A389472, COMMENTS
\cite{h474,h475,h472}.
For the first equation we establish the integral-series identity
$A(x)^2=x^2+2x^3+9Q(x)$ with $Q\in\Z[[x]]$.
A coefficient construction proves existence, integrality and uniqueness of
the normalized rational solutions of the first two equations.
Frobenius gives their congruences modulo three; binomial coefficient descent
proves both residue-class conjectures for A389472.
We also give the reversion identities, the even-power consequence
$A^{2k}\equiv x^{2k}(1+2x)^k\pmod 9$, and a factorization through
A350432 which proves the stated radius $(\sqrt5-1)/2$ for A389472
and A350432.
\end{abstract}

\section{The equations and the conjectures}

Hanna labels the square congruence and power congruence of
A389474 as \emph{conjecture} in FORMULA (3) and (4), respectively
\cite{h474}. The series is the normalized solution
$A=x+O(x^2)$ of
\begin{equation}\label{eq:A}
 A\bigl(x^2+2x^3+2A(x)^2\bigr)=3A(x)^2.
\end{equation}
Theorem~\ref{thm:nine} proves the square congruence, retaining its
two exceptional low-degree terms. The other four conjectures in the two
companion entries follow from the same integrality argument and a separate
binomial descent. Throughout, composition is formal; every inner series
has constant coefficient zero. A congruence modulo an integer means
coefficientwise divisibility in $\Z[[x]]$.

Write $B=x+O(x^2)$ for the normalized series of A389475, so that
\begin{equation}\label{eq:B}
 B(3x^2)=2x^2+B(x)^2-2B(x)^3.
\end{equation}
For A389472 write $C=\sum_{m\ge1}c_mx^m$, with the seed $c_2=1$, and
\begin{equation}\label{eq:C}
 C(x^2+x^3)=x^2(1+C(x)).
\end{equation}
This equation forces $c_1=1$ but leaves $c_2$ free. The seed is therefore
essential when identifying the OEIS sequence.
For background on integral powers and their congruences in
$\Z[[x]]$, see Heninger, Rains and Sloane \cite[Theorems 1 and 16]{hrs}.
Here the equation-specific construction supplies integrality, and an
integral unit factor supplies the extra factor nine.

\section{Existence and integrality}

\begin{lemma}\label{lem:integral}
Each of \eqref{eq:A} and \eqref{eq:B} has exactly one solution
$x+O(x^2)$ in $\Q[[x]]$. Both solutions belong to $\Z[[x]]$,
and both have coefficient of $x^2$ equal to one.
\end{lemma}
\begin{proof}
For a series $f=x+O(x^2)$ define the two residuals
\begin{align*}
 E_A(f)&=f(x^2+2x^3+2f^2)-3f^2,\\
 E_B(f)&=f(3x^2)-2x^2-f^2+2f^3.
\end{align*}
For $n\ge2$ and any rational $t$, either residual satisfies
\begin{equation}\label{eq:diagonal}
 E(f+tx^n)-E(f)=-2t x^{n+1}+O(x^{n+2}).
\end{equation}
For $E_A$, set $H_f=x^2+2x^3+2f^2$. Its change is
$H_{f+tx^n}-H_f=4t x^{n+1}+O(x^{n+2})$.
The nonlinear tail $f(y)-y$ has order at least two in $y$;
changing its inner argument affects only degrees at least $n+3$.
The changed outer coefficient contributes $tH_{f+tx^n}^{\,n}$,
whose order is at least $2n\ge n+2$.
Thus the composition changes by $4t x^{n+1}$ to the displayed precision,
whereas $-3f^2$ changes by $-6t x^{n+1}$.
For $E_B$, the substituted term has change of order $2n$ and the
cubic term has change of order $n+2$; the square term supplies
$-2t x^{n+1}$. This proves \eqref{eq:diagonal}. A perturbation
of order $n+1$ affects the residual only in degrees at least $n+2$.

If $f\in x\Z[[x]]$, all coefficients of either residual are even.
Indeed, in characteristic two each reduces to $f(x^2)-f(x)^2=0$.
This Frobenius identity follows by squaring finite polynomial truncations,
which determine each fixed coefficient of the infinite series.

Start with $f=x$, for which both residuals are $2x^3$.
Suppose the coefficients below degree $n$ have been chosen integrally,
the coefficient of $x^n$ is currently zero, and $E(f)=O(x^{n+1})$.
Let $e=\coef{n+1}{E(f)}$. It is even, so adjoining $(e/2)x^n$
annuls that coefficient by \eqref{eq:diagonal}, without changing
the lower coefficients. For $n=2$ this gives $e/2=1$.
The successive polynomials have a coefficientwise limit in $\Z[[x]]$.
Each residual coefficient depends on only finitely many coefficients,
so the limit has residual zero.

Finally, two normalized rational solutions with first difference
$t x^n$, $n\ge2$, would have residual difference with coefficient
$-2t\ne0$ in degree $n+1$ by \eqref{eq:diagonal}.
This proves uniqueness over $\Q$ and identifies every normalized rational
solution with the integral construction.
\end{proof}

\section{The congruence modulo nine and its consequences}

\begin{theorem}\label{thm:nine}
For the normalized solution of \eqref{eq:A},
\begin{equation}\label{eq:nine}
 A^2-x^2-2x^3\in9\Z[[x]].
\end{equation}
Consequently $9\mid\coef{n}{A^2}$ for every $n>3$,
as conjectured in A389474, FORMULA (3).
\end{theorem}
\begin{proof}
By Lemma~\ref{lem:integral}, write $A=x+x^2Q$ with
$Q\in\Z[[x]]$. Set $D=x^2+2x^3$ and $H=D+2A^2$.
Equation~\eqref{eq:A} becomes
\[
 H+H^2Q(H)=3A^2.
\]
The series $U=1+HQ(H)$ has constant coefficient one, so its inverse
$V=\sum_{j\ge0}(-HQ(H))^j$ belongs to $\Z[[x]]$.
Hence $H=3A^2V$. Subtracting $H=D+2A^2$ from the displayed
equation also gives $A^2-D=H^2Q(H)$. Combining the identities yields
\[
 A^2-D=9(A^2V)^2Q(H),
\]
which proves \eqref{eq:nine}.
\end{proof}

\begin{corollary}\label{cor:powers}
For every integer $k\ge0$,
\[
 A^{2k}\equiv x^{2k}(1+2x)^k\pmod9.
\]
Thus $\coef{n}{A^{2k}}$ is divisible by nine for
$n<2k$ or $n>3k$, and for $0\le j\le k$ its coefficient
in degree $2k+j$ is congruent to $2^j\binom{k}{j}$ modulo nine.
\end{corollary}
\begin{proof}
Raise \eqref{eq:nine} to the $k$th power in $(\Z/9\Z)[[x]]$
and expand $(1+2x)^k$.
\end{proof}

\begin{theorem}\label{thm:three}
For every $k\ge0$,
\begin{equation}\label{eq:frob}
 A(x)^{3k}\equiv A(x^3)^k\pmod3.
\end{equation}
The normalized solution $B$ of \eqref{eq:B} satisfies
\begin{equation}\label{eq:Bthree}
 B(x)^3\equiv B(x^3),\qquad
 B(x)^2\equiv x^2+2B(x^3)\pmod3.
\end{equation}
These prove A389474, FORMULA (4), and A389475,
FORMULA (4.a), (4.b).
\end{theorem}
\begin{proof}
For every $f\in\Z[[x]]$, reduction modulo three gives
$f(x)^3=f(x^3)$: the binomial cross terms vanish, and
$a^3=a$ for $a\in\F_3$. As before, finite truncations justify
the calculation for formal series. Taking $k$th powers gives
\eqref{eq:frob} and the first congruence of \eqref{eq:Bthree}.
Since $B(0)=0$, $B(3x^2)\equiv0\pmod3$.
Reducing \eqref{eq:B} therefore gives
$B^2\equiv x^2+2B^3\equiv x^2+2B(x^3)$.
\end{proof}

\begin{consequence}[A389475 FORMULA (3.a), (3.b)]
If $b_n=\coef{n}{B}$,
then for $n\ge1$ the odd coefficients satisfy
$\coef{2n+1}{B^2}=2\coef{2n+1}{B^3}$; for $n\ge2$,
\[
 \coef{2n}{B^2}=3^n b_n+2\coef{2n}{B^3}.
\]
\end{consequence}
\begin{proof}
Extract coefficients from \eqref{eq:B}.
In the first case $B(3x^2)$ has no odd coefficient; in the
second its degree-$2n$ coefficient is $3^n b_n$.
The exceptional term $2x^2$ contributes in neither stated range.
\end{proof}

\section{Both residue-class conjectures for A389472}

\begin{theorem}\label{thm:C}
There is exactly one solution of \eqref{eq:C} with
$c_1=c_2=1$, $C(0)=0$. It lies in $\Z[[x]]$ and satisfies
\[
 c_{2n-1}\equiv0\pmod2,\qquad
 c_{3n-1}\equiv0\pmod3\qquad(n>1).
\]
These are the two conjectures in A389472, COMMENTS.
\end{theorem}
\begin{proof}
Extracting degree $m+2$ from \eqref{eq:C} gives, for $m\ge3$,
\begin{equation}\label{eq:recurrence}
 c_m=\sum_{k=2}^{\lfloor(m+2)/2\rfloor}
 c_k\binom{k}{m+2-2k},
\end{equation}
where a binomial coefficient is zero outside $0\le j\le k$.
Every summand has $k<m$. The recurrence therefore constructs unique
integer coefficients from $c_1=c_2=1$. The coefficients below degree
five in \eqref{eq:C} hold directly for these seeds; the recurrence
establishes all subsequent coefficients.

For $p=2$ or $3$, a binomial coefficient $\binom{k}{j}$
vanishes modulo $p$ whenever $j\bmod p>k\bmod p$.
To see this, write $k=pq+r$, $0\le r<p$, and use
$(1+t)^k\equiv(1+t^p)^q(1+t)^r\pmod p$:
the possible residues of the exponent in the second factor are at most $r$.

Suppose first that $m\ge3$ is odd and argue by strong induction on $m$.
For an even $k$ in \eqref{eq:recurrence}, the lower index
$j=m+2-2k$ is odd, so the binomial coefficient is even.
For odd $k\ge3$, $c_k$ is even by induction because $k<m$.
Every summand is therefore even.

Now let $m\ge5$, $m\equiv2\pmod3$, and again use strong induction.
For a nonzero summand set $j=m+2-2k$, so
$j\equiv1+k\pmod3$.
If $k\equiv0$ or $1\pmod3$, then respectively
$j\equiv1$ or $2\pmod3$, and the binomial coefficient vanishes modulo three.
If $k\equiv2\pmod3$ and $k\ge5$, the factor $c_k$ is divisible
by three by induction. The remaining possibility $k=2$ has
$j=m-2\ge3>2$, so its binomial coefficient is zero.
All terms in \eqref{eq:recurrence} are divisible by three.
\end{proof}

The ranges cannot include $n=1$: the corresponding coefficients are
$c_1=c_2=1$. Stronger moduli, not asserted in the source, fail:
$c_3=2$ is not divisible by four, and $c_5=6$ is not divisible by nine.

\section{The companion product and the free seed}

Put $T(x)=x^2+x^3$, and let $T^{\circ j}$ denote its $j$th
compositional iterate, with $T^{\circ0}(x)=x$.

\begin{theorem}\label{thm:product}
The normalized series $D=x+O(x^2)$ of A350432 exists uniquely and is
\[
 D(T(x))=xD(x),\qquad
 D(x)=x\prod_{j\ge0}\bigl(1+T^{\circ j}(x)\bigr).
\]
It has nonnegative integer coefficients. For every $s\in\Q$, the
unique solution of \eqref{eq:C} with $C(0)=0$, $c_1=1$, $c_2=s$ is
\begin{equation}\label{eq:seed}
 C_s(x)=x+sD(x)^2.
\end{equation}
In particular $C=C_1$. Both $D$ and $C_s$, for $s\ne0$, have radius of
convergence $\rho=(\sqrt5-1)/2$; $C_0=x$ is a polynomial.
\end{theorem}
\begin{proof}
Write $D=xE$. The equation for $D$ is equivalent to
$E(x)=(1+x)E(T(x))$, $E(0)=1$. Iterating gives
\[
 E(x)=\prod_{j=0}^{N-1}(1+T^{\circ j}(x))\,E(T^{\circ N}(x)).
\]
Since $\ord(T^{\circ N})=2^N$, the last factor tends coefficientwise
to one. This proves uniqueness and the product formula. Conversely the
product is a well-defined integral formal series, with nonnegative
coefficients, and shifting its factors verifies the equation.
Setting $F_j=1+T^{\circ(j-1)}$ gives
$F_1=1+x$ and $F_{j+1}=1+(F_j-1)^2F_j$, exactly the product
in A350432, FORMULA (3) \cite{h350}.

The identity $D(T)^2=x^2D^2$ shows directly that $x+sD^2$ solves
\eqref{eq:C} and has the prescribed seed. The triangular recurrence
\eqref{eq:recurrence}, with $c_2=s$, proves uniqueness over $\Q$.

For completeness, the same product proves the radius asserted in both
entries \cite{h472,h350}. If $0<r<\rho$, then $T(r)<r$ and its
iterates $r_j=T^{\circ j}(r)$ decrease to zero: a positive limit $\ell$
would solve $\ell+\ell^2=1$, hence would be $\rho>r$.
Eventually $r_j\le1/4$, so $r_{j+1}=r_j^2(1+r_j)\le r_j/2$.
Thus $\sum_j r_j<\infty$. On $|z|\le r$ we have
$|T^{\circ j}(z)|\le r_j$, so the product converges uniformly there
to an analytic function. Its Taylor coefficients are the formal
coefficients above. This proves a radius at least $\rho$.
For any fixed $N$, positivity gives, on $0<r<\rho$,
\[
 D(r)\ge r\prod_{j=0}^{N-1}(1+T^{\circ j}(r)).
\]
As $r\uparrow\rho$, the right side tends to
$\rho(1+\rho)^N$. Letting $N$ grow shows $D(r)\to+\infty$.
Consequently $D$ cannot be analytic across $\rho$ and its radius is
exactly $\rho$. Equation~\eqref{eq:seed} gives the same divergence
and unboundedness, hence the same radius, for every nonzero rational $s$.
\end{proof}

\begin{corollary}\label{cor:seeds}
For every integer seed $s$, the series $C_s$ is integral and satisfies
$2\mid[x^{2n-1}]C_s$ and $3\mid[x^{3n-1}]C_s$ for $n>1$.
If $V=D^{\langle-1\rangle}$, then
$V(xD(x))=x^2+x^3$, as in A350432, FORMULA (2).
\end{corollary}
\begin{proof}
For degrees at least two, \eqref{eq:seed} multiplies the coefficients
of $C_1$ by $s$, so Theorem~\ref{thm:C} proves the congruences.
Apply $V$ to $D(T)=xD$ for the reversion identity.
\end{proof}
The all-seed statement and the factorization \eqref{eq:seed} are
extensions not posed as conjectures in the cited entries.

\section{Reversion, sharpness and source coverage}

Since each series has linear coefficient one, its compositional inverse
exists over $\Z[[x]]$. Indeed, the coefficient of degree $n$ in
$F(R(x))=x$, with $F=x+O(x^2)$, determines $[x^n]R$
with coefficient one from already known integer coefficients.
The resulting right inverse is also a left inverse: composition with
$F=x+O(x^2)$ preserves the first nonzero coefficient of a difference,
and is therefore injective.

\begin{consequence}[Reversion and iteration]
Let $R=A^{\langle-1\rangle}$, $S=B^{\langle-1\rangle}$ and
$W=C^{\langle-1\rangle}$. Then
\begin{align*}
 A(2x^2+R^2+2R^3)&=3x^2, & R(3x^2)&=2x^2+R^2+2R^3,\\
 B(3S^2)&=2S^2+x^2-2x^3, & C(W^2+W^3)&=(1+x)W^2.
\end{align*}
These are A389474 (2), A389475 (2), and A389472 (2),
including the reversion identity in A389474 EXAMPLE.
Also A389472 (3) holds as the Laurent-series identity
\[
 C(x)=\frac{C(x^4+2x^5+2x^6+3x^7+3x^8+x^9)}{(x^3+x^4)^2}
       -\frac{1+x^2}{x^2}.
\]
\end{consequence}
\begin{proof}
Substitute the compositional inverses into the defining equations.
For $T=x^2+x^3$, composing \eqref{eq:C} with $T$ gives
$C(T(T))=T^2(1+x^2(1+C(x)))$.
Expanding $T(T)$ and rearranging proves the Laurent-series formula;
the apparent negative powers cancel by this identity.
\end{proof}

\begin{corollary}\label{cor:273}
The normalized series $P$ of A273095,
$P(x^2)=P(x)^2-2P(x)^3$, exists uniquely in $\Z[[x]]$.
If $J=P^{\langle-1\rangle}$ and $L(x)=x^2-2x^3$, then
\[
 P(J(x)^2)=L(x),\qquad J(L(x))=J(x)^2,\qquad
 P(J(x)^{2^r})=L^{\circ r}(x)\quad(r\ge0).
\]
These include the FORMULA and iterated EXAMPLE identities of A273095
\cite{h273}.
\end{corollary}
\begin{proof}
For $E_P(f)=f(x^2)-f^2+2f^3$, the coefficient diagonal is again
$-2t x^{n+1}$ as in \eqref{eq:diagonal}. Reduction modulo two makes
the residual zero, and $E_P(x)=2x^3$. The construction and uniqueness
argument of Lemma~\ref{lem:integral} therefore apply unchanged.
The equation $P(x^2)=L(P(x))$ iterates to
$P(x^{2^r})=L^{\circ r}(P(x))$. Substitute $x=J(y)$.
The case $r=1$, followed by composition with $J$, gives the first
two identities.
\end{proof}

The following stronger moduli are not asserted in the source;
they fail for the displayed coefficients.
The initial coefficients $A=x+x^2+4x^3+O(x^4)$ and
$B=x+x^2+7x^3+O(x^4)$ give
\[
 \coef{4}{A^2}=9,\qquad
 \coef{4}{B^3-B(x^3)}=3,\qquad
 \coef{4}{B^2-x^2-2B(x^3)}=15.
\]
Thus the moduli nine, three and three in the respective statements
cannot be replaced by twenty-seven, nine and nine.

\begin{consequence}[A389474 COMMENTS: Catalan comparison]
For $K=x+K^2$, the unique normalized Catalan series,
\[
 x^2+2xK^2=K^2-K^4,\qquad K(x^2+2xK^2)=K^2.
\]
\end{consequence}
\begin{proof}
Here $K$ is $x$ times the ordinary Catalan generating function
\cite{catalan}. Its uniqueness follows recursively from $K=x+K^2$.
Substitute $x=K-K^2$ to obtain the first identity; applying $K$ to
the identity $K^2-K^4=(z-z^2)|_{z=K^2}$ proves the second, since
$z-z^2$ is the compositional inverse of $K$.
\end{proof}

\begingroup\small\raggedright
\begin{longtable}{@{}>{\raggedright\arraybackslash}p{0.39\linewidth}>{\raggedright\arraybackslash}p{0.55\linewidth}@{}}
\toprule
Source item & Coverage\\
\midrule
A389474 FORMULA (3): $9\mid[x^n]A^2$, $n>3$ (``conjecture'') &
Proved: Theorem~\ref{thm:nine}.\\
A389474 FORMULA (4): $3\mid[x^n](A^{3k}-A(x^3)^k)$,
$n,k\ge1$ (``conjecture'') & Proved: Theorem~\ref{thm:three}.\\
A389475 FORMULA (4.a): $B^3\equiv B(x^3)\pmod3$ (``conjecture'') &
Proved: Theorem~\ref{thm:three}.\\
A389475 FORMULA (4.b): $B^2\equiv x^2+2B(x^3)\pmod3$
(``conjecture'') & Proved: Theorem~\ref{thm:three}.\\
A389472 COMMENTS: ``Conjecture: a(2*n-1) = 0 (mod 2) for n > 1.'' &
Proved: Theorem~\ref{thm:C}, with $c_2=1$.\\
A389472 COMMENTS: ``Conjecture: a(3*n-1) = 0 (mod 3) for n > 1.'' &
Proved: Theorem~\ref{thm:C}, with $c_2=1$.\\
A389474 (1), A389475 (1) & Proved existence and integrality:
Lemma~\ref{lem:integral}.\\
A389474 (2), reversion EXAMPLE & Proved by substitution and reversion.\\
A389475 (2), (3.a), (3.b) & Proved by substitution and coefficient extraction.\\
A389472 (1) & Proved seeded existence and uniqueness:
Theorem~\ref{thm:C}.\\
A389472 (2), (3) & Proved by reversion and two compositions.\\
A389474 Catalan comparison & If $K=x+K^2$, then
$x^2+2xK^2=K^2-K^4$ and $K(x^2+2xK^2)=K^2$;
the normalized solution $K$ is unique.\\
A389472 and A350432 COMMENTS: radius $(\sqrt5-1)/2$ &
Proved: Theorem~\ref{thm:product}.\\
A350432 FORMULA (1)--(3) and generating-method recurrence &
Proved: Theorem~\ref{thm:product} and Corollary~\ref{cor:seeds}.\\
A273095 defining equation, FORMULA reversion and iterated EXAMPLE &
Proved: Corollary~\ref{cor:273}.\\
All integer seeds in A389472; $C_s=x+sD^2$ &
Not posed; proved: Theorem~\ref{thm:product} and
Corollary~\ref{cor:seeds}.\\
A273095 asymptotic constants; displayed decimal evaluations &
Not covered.\\
\bottomrule
\end{longtable}
\endgroup

\section{Formal statements and finite checks}

The following pinned Lean files contain declarations with the specified
mathematical scope. The namespaces share the prefix
\nolinkurl{D5.S1.Recurrence.}; the table gives the remaining module and
declaration names. Each filename links directly to its source at commit
\nolinkurl{074df09ddbbef1843ddb383577bfaa827f753e36}.

\begingroup\small
\begin{longtable}{@{}p{0.55\linewidth}p{0.39\linewidth}@{}}
\toprule
Module and declaration & Statement\\
\midrule
\endhead
\href{\leanroot/D5/S1/Recurrence/Invariants/HannaCompositionSquareModNine.lean}
{\nolinkurl{Invariants.HannaCompositionSquareModNine}}\newline
\nolinkurl{result} & Existence and divisibility by nine for every
normalized rational solution, $n>3$.\\[3pt]
Same file: \nolinkurl{FormulaFour.resultFour} & A389474 (4),
all $n,k\ge1$, with integer multiples of three.\\[3pt]
\href{\leanroot/D5/S1/Recurrence/Invariants/HannaTripleSquareFrobeniusModThree.lean}
{\nolinkurl{Invariants.HannaTripleSquareFrobeniusModThree}}\newline
\nolinkurl{result}, \nolinkurl{FormulaFourB.resultFourB} &
A389475 (4.a), (4.b), respectively, for every normalized rational solution
and every coefficient index.\\[3pt]
Same file: \nolinkurl{integral_frobenius} &
$3\mid[x^n](f^{3k}-f(x^3)^k)$ for every integral series and all $n,k\ge0$.\\[3pt]
\href{\leanroot/D5/S1/Recurrence/Parity/SquareCubeSubstitutionThreeIndexModThree.lean}
{\nolinkurl{Parity.SquareCubeSubstitutionThreeIndexModThree}}\newline
\nolinkurl{result} & $3\mid c_{3n-1}$ for $n>1$, with seed $c_2=1$.\\[3pt]
\href{\leanroot/D5/S1/Recurrence/Parity/SquareCubeSubstitutionOddIndexParity.lean}
{\nolinkurl{Parity.SquareCubeSubstitutionOddIndexParity}}\newline
\nolinkurl{hanna_conjecture}, \nolinkurl{generating_equation},
\nolinkurl{generating_unique} &
$2\mid c_{2n-1}$ for $n>1$; the integral recurrence is identified with
the normalized, seeded solution of \eqref{eq:C}.\\
\bottomrule
\end{longtable}
\endgroup

The proofs in this note establish all six conjectures without a finite
cutoff. Corollary~\ref{cor:powers}, the reversion formulas, the iterated
identities, the product, all-seed factorization and radius results are
written proofs. The four linked files cover the six congruence
conjectures, not these additional statements.
The accompanying exact-integer checker tests the nonlinear equations through
degree 97, the linear equation and congruences through degree 96, and matches
the 24, 23 and 41 displayed source coefficients.
It checks exponents $0\le k\le33$ in the Frobenius identity and
$0\le k\le49$ in Corollary~\ref{cor:powers}; higher exponents have
zero coefficients in this truncation by order.
It also matches the 34 and 43 displayed coefficients of A273095 and
A350432, respectively, and checks their equations, the product,
the factorization, and reversion and iterated identities through the
stated finite degree. The radius assertions follow from the analytic
proof, not a finite check.

\begingroup\small
\begin{thebibliography}{9}
\bibitem{hrs} Nadia Heninger, E. M. Rains and N. J. A. Sloane,
\emph{On the integrality of $n$th roots of generating functions},
arXiv:math/0509316v4, April 8, 2006, Theorems 1 and 16.
\url{https://arxiv.org/abs/math/0509316v4}.
\bibitem{h474} Paul D. Hanna, \emph{A389474}, The On-Line Encyclopedia of
Integer Sequences, author date October 24, 2025; entry revision 11,
October 26, 2025. FORMULA (1)--(4), COMMENTS and EXAMPLE.
\url{https://oeis.org/A389474}.
\bibitem{h475} Paul D. Hanna, \emph{A389475}, The On-Line Encyclopedia of
Integer Sequences, author date October 26, 2025; entry revision 8,
October 27, 2025. FORMULA (1)--(4.b).
\url{https://oeis.org/A389475}.
\bibitem{h472} Paul D. Hanna, \emph{A389472}, The On-Line Encyclopedia of
Integer Sequences, author date October 15, 2025; entry revision 7,
October 17, 2025. COMMENTS and FORMULA (1)--(3).
\url{https://oeis.org/A389472}.
\bibitem{h273} Paul D. Hanna, \emph{A273095}, The On-Line Encyclopedia of
Integer Sequences, May 15, 2016; entry revision 18, September 2, 2017.
\url{https://oeis.org/A273095}.
\bibitem{h350} Paul D. Hanna, \emph{A350432}, The On-Line Encyclopedia of
Integer Sequences, December 30, 2021; entry revision 13, January 26, 2024.
\url{https://oeis.org/A350432}.
\bibitem{catalan} \emph{A000108: Catalan numbers}, The On-Line Encyclopedia
of Integer Sequences, FORMULA: $C(x)=1+xC(x)^2$.
\url{https://oeis.org/A000108}.
\end{thebibliography}
\endgroup

\noindent\small\textcopyright\ 2026 Haobo Ma and Wenlin Zhang.
Licensed under Creative Commons Attribution 4.0 International
(CC BY 4.0; \texttt{cc-by-4.0}).
\end{document}
